\subsection{Laplace-Type Operators and Asymptotic Behavior}
  \label{subsec:laplace-type-operators}

  We consider an elliptic operator of Laplace type in three spatial dimensions,
  \begin{equation}
    A = -\nabla^2 + V(x),
  \end{equation}
  acting on a bounded domain $\Omega \subset \mathbb{R}^3$
  with boundary conditions ensuring self-adjointness and positivity.

  In the weak-field regime and away from boundaries,
  the leading-order behavior of the Green kernel is governed
  by the principal symbol of the operator.
  Neglecting lower-order contributions, the Green function satisfies
  \begin{equation}
    -\nabla^2 G(x,y) = \delta(x-y).
  \end{equation}

  In three dimensions, the fundamental solution of the Laplacian
  is well known:
  \begin{equation}
    G(x,y) = \frac{1}{4\pi |x-y|}.
  \end{equation}
  The relation between resolvents, heat kernels, and spectral
  invariants has been extensively studied in the mathematical
  physics literature~\cite{gilkey1995,vassilevich2003}.

  More generally, for operators of the form $A = -\nabla^2 + V(x)$
  with sufficiently regular and bounded $V$,
  the Green kernel admits the asymptotic expansion
  \begin{equation}
    G(x,y) \sim \frac{1}{4\pi |x-y|}
    \quad \text{as } |x-y| \to \infty,
  \end{equation}
  up to exponentially suppressed or subleading corrections,
  provided no mass gap is introduced.
